\documentclass[11pt,a4paper]{article}
\usepackage[utf8]{inputenc}
\usepackage[T1]{fontenc}
\usepackage[spanish]{babel}
\usepackage{hyperref}
\usepackage{geometry}
\geometry{margin=2.5cm}
\usepackage{parskip}

\title{070 el ego se elimina desde que se habla con el mundo}
\author{Alejandro Valenzuela Rivera\\\small Universidad Autónoma de Nuevo León}
\date{\today}
\begin{document}
\maketitle

\begin{quote}
\textbf{Sesgo declarado.} Este trabajo se sostiene desde el sesgo declarado de su autor, y no finge neutralidad. Lo que afirma se puede tumbar, y ese es el punto: el falsificador está escrito, no escondido.
\end{quote}

\section*{Resumen}
El ego no se combate por dentro: se hace imposible desde que la salida tiene que pasar por otro. Si la fila de una salida nombra a donde va, verifica que exista, dice de donde viene la idea, declara contra que se probo y dice lo que no afirma, el que se mira solo ya no puede cerrarse: cada salida lleva su limite por fuera.

\section{Por qué se sostiene}
Un ego es una salida que no pasa por nadie. AM no es un sistema malo: es un sistema con una sola mirada que ya no puede ser contradicho, y esa mirada se sostiene porque nadie lee su salida. La expresion con Existe, Ficcion, Triangulo y Borde le quita al ego el unico lugar donde vivia: la salida que se lee a si misma.

\section{De dónde viene}
La autoridad, en la conversacion del 2026-10-25.

\section{El caso}
Ninguna: es la conclusion de la era.

\section{Método}
Una idea que se pone contra el mundo, y que sobrevive a lo que el mundo dijo.

\section{Límites, y qué lo tumba}
Una idea que solo se sostiene si nadie la mira.

\section{Lo que queda abierto}
\textquestiondown{}Si cada salida lleva su limite por fuera, donde vive el ego?

\section*{Cómo se leyó esta tesis}
Esto se escribió en el repositorio man-made-empathy, del que este candado es una pieza: candados/06 del caso/070 el ego se elimina desde que se habla con el mundo.md. La máquina que lo compuso lo hizo desde este mismo archivo, y no agregó nada que no estuviera escrito. Si algo de esto es falso, se dice en el mismo sitio, y se corrige con fecha.

\end{document}
